% Current implementation status; evidence lives in docs/phase4_* and docs/internal/.
\begin{frame}{Phase 4：残時間推定と容量調停の進捗（2026-10-15）}
  \begin{columns}[T,onlytextwidth]
    \begin{column}{.48\textwidth}
      \textbf{残りの実行形状を実測で価格付け}\par
      \small
      \begin{itemize}
        \item 残りepochのexecutor呼出しと公開処理を数える。
        \item 実行中の呼出しは経過時間に応じた残量で推定。
        \item 角度poolingも実submit順でheld-out認証。
        \item 未認証classがあればETAはunknownを返す。
      \end{itemize}
    \end{column}
    \begin{column}{.48\textwidth}
      \textbf{同じvolumeの独立ledgerを調停}\par
      \small
      \begin{itemize}
        \item 通常runtimeからcoordinatorを明示指定。
        \item fork・勾配の子も束縛を継承。
        \item 再開時の設定不一致を全rankで拒否。
        \item 専用rootの増分で外部jobの宣言超過を検出。
      \end{itemize}
    \end{column}
  \end{columns}
  \smallnote{短いrunの48地点中30地点で回答、中央値誤差8.6\%、区間coverage 83.3\%。回答率改善の比較実証は残る。}
  \smallnote{Linux回帰：buffered 105成功、O\_DIRECT 112成功・2 skip（Mac専用）。実MPI 1/2 rankと回収を含む。}
  \takeaway{P4-001はDONE。\textbf{P4-002・009とPhase 4全体はIN\_PROGRESS}。複数host調停・全writerの故障回収・電源断modelは残る。}
  \note{根拠：docs/phase4\_handoff\_2026-10-15.md、docs/phase4\_exact\_future\_eta\_implementation.md、docs/phase4\_disk\_coordination\_implementation.md。容量調停はnode-localのopt-in。共有filesystemは拒否し、strict\_capacity\_guaranteeは立てない。layout migrationのtransient・Linux VERIFY・O\_DIRECT予測、tile/cache/gradient自動運用と最終比較matrixもbacklogへ残す。}
\end{frame}

\begin{frame}{Phase 4：回収・移行判断・比較証跡の補強}
  \small
  \begin{tabularx}{\textwidth}{@{}p{.18\textwidth}XX@{}}
    \toprule 対象 & 追加した確認 & 未達として残す条件 \\
    \midrule
    Catalog予約 & DBのpath・device・inode・birthtimeへ帰属。DB／journalの観測を回収計画とreceiptに保存。 & 旧token、birthtime不明、作成時のidentity未確定は回収を拒否。 \\
    \addlinespace
    移行のbaseline & prefixのcoverage・drift・executor classを確認。現在の証拠が無効なら古いestimateも失効。 & 不安定なbaselineはunknown。評価用の認証をproductionへ転用しない。 \\
    \addlinespace
    比較matrix & 8 groupの代表case、2 variant・warmup・paired trialを検証。大状態とcross-hostは別slot。 & 未実行・予算不足・unsupported・数値不良を残す。全axisの受入は別判定。 \\
    \bottomrule
  \end{tabularx}
  \medskip
  \smallnote{統合高速1081成功・36 skip。Linux buffered／O\_DIRECTは各196成功。実statx birthtimeと別volume拒否も確認。}
  \takeaway{不確かな証拠を完了扱いせず、\textbf{拒否理由と不足する検証を記録する}。P4-003・004・005・009はIN\_PROGRESS。}
  \note{根拠：docs/phase4\_uncommitted\_recovery\_implementation.md、docs/phase4\_layout\_migration\_implementation.md、docs/phase4\_comparison\_matrix.md。birthtimeはfilesystemの分解能に依存する。比較matrixはrepresentative coverage onlyで、Phase 4 acceptanceはnot-assessed。損益分岐規則は変更しない。}
\end{frame}

\begin{frame}{Phase 4：実行経路とクラッシュ境界の追加検証（2026-10-16）}
  \begin{columns}[T,onlytextwidth]
    \begin{column}{.48\textwidth}
      \textbf{P4-002：実行経路を明示}\par
      \small
      \begin{itemize}
        \item 実dispatchと共通のrouteから有界stage DAGを生成。
        \item 読み書き・kernel・通信・公開の依存を保持。
        \item inclusive親時間は一度だけ計上。
        \item 件数・MPI payload・kernel・storage coverageを分離。
      \end{itemize}
    \end{column}
    \begin{column}{.48\textwidth}
      \textbf{P4-009：起動・export・電源断模型}\par
      \small
      \begin{itemize}
        \item catalog作成前intentとinode束縛を段階的に永続化。
        \item payload／chunkとexportの故障後予約を回収。
        \item export先もvolume全体の使用量へ反映。
        \item SQLiteをEXTRAへ変更しjournal削除を同期。
      \end{itemize}
    \end{column}
  \end{columns}
  \smallnote{統合高速1317成功・38 skip。実2 host・4 rankで調停の8 case成功（buffered／O\_DIRECT）。}
  \takeaway{P4-002・009はIN\_PROGRESS。Linux最終fault matrix、監視費用・endpointのheld-out評価を継続する。}
  \note{根拠：docs/phase4\_executor\_stage\_dag\_implementation.md、docs/phase4\_catalog\_startup\_recovery\_implementation.md、docs/phase4\_sqlite\_power\_loss\_implementation.md。電源断模型は成功したVFS syncのみをdurable imageへ反映する。実装済みDBとjournalの模型であり、実hardwareやtorn writeの保証は主張しない。payload/storage/Frameの未確定量はunknownとし、logical state bytesで代用しない。}
\end{frame}
